\documentclass[a4paper, 12pt, pdftex]{article}
\usepackage{graphicx}
%\usepackage[pdftex, colorlinks]{hyperref}
\usepackage{amsmath, amsthm}
\usepackage{enumitem}
\usepackage{hyperref}

\usepackage[T1]{fontenc}
\usepackage[full]{textcomp}

%\usepackage{arev}
%\usepackage{fourier}
%\usepackage{fouriernc}
%\usepackage[utopia]{mathdesign}
%\usepackage[garamond]{mathdesign}
\usepackage[charter]{mathdesign}
%\usepackage{kpfonts}

%\usepackage{times}           % Times + Helvetica + Courier + Symbol 
%\usepackage{mathptmx}        % Times + Helvetica + Courier + Times Math
%\usepackage{mathpazo}        % Palatino + Helvetica + Courier + Palatino 
%\usepackage{newcent}         % New Century Schoolbook + Avant Garde + Courier
%\usepackage{bookman}         % Bookman + Avant Garde + Courier

%\DeclareGraphicsExtensions{.eps , .pdf}

\usepackage{fullpage}
%\setlength{\topmargin}{0 cm}
%\setlength{\headheight}{-2.5 cm}
%\setlength{\headsep}{0 cm}
%\advance\textheight by 4.0 cm
%\advance\textwidth by 2.0 cm
%\advance\oddsidemargin by -1.0 cm
%\advance\evensidemargin by -1.0 cm

\usepackage{amsmath}
\usepackage[utf8]{inputenc}
%\usepackage[noend]{algorithm, algpseudocode}

\setlength{\parindent}{0cm}
\sloppy
\pagestyle{plain}

\theoremstyle{plain}
\newtheorem{theorem}{Theorem}
\newtheorem{corollary}{Corollary}
\newtheorem{lemma}{Lemma}
\newtheorem{conjecture}{Conjecture}

\theoremstyle{definition}
\newtheorem{examples}{Examples}
\newtheorem{definition}{Definition}

\DeclareMathOperator{\hash}{hash}

%\date{}




\begin{document}
%\thispagestyle{empty}

\begin{center}
\huge Annie
\end{center}
\begin{center}
\large Montgomery McLaughlin Gallot Woltman algorithm
\end{center}
\begin{center}
\it \large Yves Gallot
\end{center}

\bigskip
\bigskip

Notation:\\
"$\bmod\, M$" is the operator: $y = x \mod M \Rightarrow y \in [0, M - 1]$.\\
The congruence relation is noted $a \equiv b \pmod M$.

\subsection*{Montgomery modular multiplication}
Peter~L.~Montgomery presented a method for multiplying two integers modulo $N$ while avoiding division by $N$ \cite{Montgomery1}.
\begin{theorem}
Let $N$ be a positive integer, $R > N$, $\gcd(R,\, N) = 1$, $0 \leq x, y < N$. If
\begin{enumerate}[label={(\arabic*)}]
\item $T = x\,y$,
\item $m = -T\,N^{-1} \mod R$,
\item $r = (T + m\, N)\, /\, R$,
\item $t = r - N$ if $r \geq N$ and $t = r$ otherwise
\end{enumerate}
then $t = R^{-1}\, x\,y \mod N$.
\end{theorem}
\begin{proof}
(2) $\Rightarrow m\, N \equiv -T \pmod R$ then $r$ is an integer.
(3) $\Rightarrow r \equiv T\, R^{-1} \pmod N$.\\
$0 \leq m\, N < R\, N$ and $0 \leq T < N^2 < R\, N$ then $0 \leq r < 2\,N$.
\end{proof}

In practice, $x$ is converted into Montgomery form $x^* = R\, x \mod N$
($y$, $z$ resp.).\\ We have
$z = x\,y\mod N \Leftrightarrow z^*= R^{-1}\, x^*\,y^* \mod N$.
Conversion out of Montgomery form is done by computing $x = R^{-1}\, x^* \mod N$.

\subsection*{Montgomery modular multiplication without test}

The test (4) can be removed if $R > 4\,N$ and $0 \leq x, y, r < 2\,N$.
\begin{theorem}
Let $N$ be a positive integer, $R > 4\,N$, $\gcd(R,\, N) = 1$, $0 \leq x, y < 2\,N$. If
\begin{enumerate}[label={(\arabic*)}]
\item $T = x\,y$,
\item $m = -T\,N^{-1} \mod R$,
\item $r = (T + m\, N)\, /\, R$
\end{enumerate}
then $r \equiv R^{-1}\, x\,y \pmod N$ and $0 \leq r < 2\,N$.
\end{theorem}
\begin{proof}
$0 \leq m\, N < R\, N$ and $0 \leq T < 4\,N^2 < R\, N$ then $0 \leq r < 2\,N$.
\end{proof}

\subsection*{McLaughlin modular multiplication (without test)}

Philip~B.~McLaughlin~Jr. extended the method using a second integer Q \cite{McLaughlin1}.

\begin{theorem}
Let $N$ be a positive integer, $R > 4\,N$, $Q > R$,
$\gcd(R,\, N) = \gcd(Q,\, R) = 1$, $0 \leq x, y < 2\,N$. If
\begin{enumerate}[label={(\arabic*)}]
\item $T_R = x\,y \mod R$, $T_Q = x\,y \mod Q$,
\item $m = -T_R\,N^{-1} \mod R$,
\item $r = (T_Q + m\, N)\, R^{-1} \mod Q$
\end{enumerate}
then $r \equiv R^{-1}\, x\,y \pmod N$ and $0 \leq r < 2\,N$.
\end{theorem}
\begin{proof}
$Q > 2\,N$ and $r$ in $[0, Q - 1]$ is unique. Then $r$ is the solution.
\end{proof}

\subsection*{Gallot \& Woltman method}

The following method has emerged from a discussion on the mersenneforum \cite{Gallot1}.\\

$Q = R + 2$ is fixed. $R$ must be odd such that $\gcd(Q,\, R) = 1$.\\
Because $R^{-1} = -2^{-1} \mod Q$, we define the new form
$$x^* = 2^{-1}\,R\, x \mod N$$
($y$, $z$ resp.). Now, we have
$z = x\,y\mod N \Leftrightarrow z^*= 2\,R^{-1}\, x^*\,y^* \mod N$.\\
Note that $N$ must be odd such that $2^{-1} \mod N$ exists.
If $T_R \neq 0$ and $m' = T_R\,N^{-1} \mod R$ then $m = R - m' = -(m' + 2) \mod Q$. If $T_R = 0$ then $T_Q = m' = 0$ and $r = N$.\\
Finally, if $0 \leq x < 2\,N$ then $0 \leq r \leq 4\,N$. New conditions are
needed:
\begin{theorem}
Let $N$ and $R$ be positive odd integers, $\gcd(R,\, N) = 1$, $R > 16\,N$,
$Q = R + 2$, $0 \leq x, y < 4\,N$. If
\begin{enumerate}[label={(\arabic*)}]
\item $T_R = x\,y \mod R$, $T_Q = x\,y \mod Q$,
\item $m' = T_R\,N^{-1} \mod R$,
\item $r' = \left((m' + 2)\, N - T_Q\right) \mod Q$
\end{enumerate}
then $r' \equiv 2\,R^{-1}\,x\,y \pmod N$ and $0 \leq r' < 4\,N$.
\end{theorem}
\begin{proof}
We have $r' = 2\,r = 2\,(T + m\, N)\, /\, R$.
$0 \leq m\, N < R\, N$ and $0 \leq T < 16\,N^2 < R\, N$ then $0 \leq r' < 4\,N$.
\end{proof}

\medskip

The implementation in Annie uses $R = 2^n - 1$. The polynomial multiplication modulo $R$ is a cyclic  (positive-wrapped) convolution and modulo $Q$ is a negacyclic (negative-wrapped) convolution.\\
Number Theoretic Transforms achieve $O(n \log n)$ time complexity (they are a divide and conquer technique).\\
The code for the cyclic and negacyclic convolutions is shared. The twiddle factors are the $n$\textsuperscript{th} root of unity for $R$ and the $n$\textsuperscript{th} root of $-1$ for $Q$.


\begin{thebibliography}{99}

\bibitem{Gallot1} Yves~Gallot, George~Woltman \textit{et~al}., \emph{A new algorithm for generic modular reduction}, mersenneforum.org, Aug.--Sep.~2023,\\ \url{https://www.mersenneforum.org/node/22013}.

\bibitem{McLaughlin1} Philip~B.~McLaughlin~Jr., \emph{New frameworks for Montgomery’s modular multiplication method}, Math. Comp. \textbf{73} (2004), 899--906,\\ DOI: \url{https://doi.org/10.1090/S0025-5718-03-01543-6}.

\bibitem{Montgomery1} Peter~L.~Montgomery, \emph{Modular multiplication without trial division}, Math. Comp. \textbf{44} (1985), 519--521,\\ DOI: \url{https://doi.org/10.1090/S0025-5718-1985-0777282-X}.


\end{thebibliography}

\end{document}

